\documentclass[../mainnull]{subfiles}
\setcounter{chapter}{9}
\begin{document}

\chapter{Deslizamiento y sincronismo}
\begin{center}
\img{images/08/logo.jpg}{Motor Trifásico}
\end{center}

\subsection{Sincronismo}

La velocidad de giro del campo magnético se conoce como velocidad de sincronismo $n_s$. El motor síncrono gira a la misma velocidad que el campo magnético. 

\begin{equation}
    N=f \frac{60}{p}
\end{equation}

donde:

$N$: Numero de revoluciones por minuto.

$f$:Frecuencia de la red en Hertz.

$p$: Número de par de polos del motor.

\begin{table}[h!]
    \centering
    \begin{tabular}{|c|c|c|}\hline
        $N^o$ de Polos &  50 Hz & 60 Hz \\ \hline \hline
        2 polos & 3000 &3600 \\ 
        4 polos & 1500 & 1800 \\
        6 polos & 1000 & 1200 \\
        8 polos & 750 & 900 \\
        12 polos & 500 & 600 \\
        16 polos & 375 & 450 \\
        24 polos & 250 & 300 \\ \hline
    \end{tabular}
    \caption{Velocidad Sincrónica para motores a 50 y 60 Hz.}
    \label{tab:my_tableherz}
\end{table}

\subsection{Deslizamiento}

\web{https://www.youtube.com/watch?v=duRA-drcfGQ}{Principio de funcionamiento de motor asincrónico}

\info{Definición de deslizamiento}{El \emph{deslizamiento o resbalamiento} en una máquina eléctrica es la diferencia relativa entre la velocidad del campo magnético (velocidad de sincronismo) y la velocidad del rotor. En motores asincronicos esta velocidad no es igual.}
Las siguientes expresiones son equivalentes para hallar el deslizamiento:
\begin{equation}
    S=\frac{\omega_s - \omega_m}{\omega_s} \cdot 100\%= \frac{n_s-n_m}{n_s}\cdot 100\%
\end{equation}

Donde:

$S$: Velocidad de deslizamiento (expresada con base por unidad o en porcentaje).

$\omega_{s}$: Velocidad angular de sincronismo en radianes por segundo.

$\omega_{m}$: Velocidad angular del rotor en radianes por segundo.

$n_{s}$: Velocidad angular de sincronismo en revoluciones por minuto.

$n_{m}$: Velocidad angular del rotor en revoluciones por minuto.

El deslizamiento es especialmente útil cuando analizamos el funcionamiento del Motor asíncrono ya que estas velocidades son distintas. El voltaje inducido en el bobinado rotórico de un motor de inducción depende de la velocidad relativa del rotor con relación a los campos magnéticos.
Es posible expresar la velocidad mecánica del eje del rotor, en términos de la velocidad de sincronismo (velocidad del campo magnético) y de deslizamiento.
\begin{equation}
n_{m}=n_{e}(1-s)
\end{equation}
\info{Rango de deslizamiento}{
En el Motor asíncrono $0<s<1$ ($-1<s<0$ funcionando como generador)(Depende la aplicación, alrededor de 5\%) . En el Motor síncrono $S=0$
}

\subsection{Par Motor y velocidad}

\web{https://www.youtube.com/watch?v=vcr5pkT6p1o}{Arranque motores de inducción}


\img{images/08/curva}{Curva Parmotor-velocidad}

La velocidad nominal siempre será inferior a la velocidad de sincronismo o velocidad del campo magnético del estátor, debido al deslizamiento.

El par resistente aumenta a medida que aumenta la velocidad y el par motor disminuye al aumentarla velocidad.

Alrededor de la velocidad nominal existe una zona estable  en la que el motor puede oscilar en función del par resistente,  variando ligeramente la velocidad. Si lo sacamos de esa zona el motor saldrá de sincronismo y se parará.

La Intensidad de arranque es  muy elevada del orden de 8 veces superior a la Intensidad nominal.

De ahí la necesidad de utilizar sistemas de arranque que minimicen la Intensidad de arranque.

\begin{equation}
    P=\omega \cdot M
\end{equation}

Donde:

$P$ : Potencia en Watts.

$\omega$: velocidad del motor en (radianes/segundo)

$M$:  Par motor en Newton/metro.

Multiplicaremos por $\frac{2\cdot\pi}{60}$ para pasar de R.P.M. a radianes/segundo.

\actividad{Resolver}
\begin{enumerate}
    \item Un motor asíncrono trifásico de rotor en cortocircuito posee una velocidad de campo $n_S$ de 3000 R.P.M ¿Cuál será el deslizamiento de rotor a plena carga si se mide con un tacómetro una velocidad de 2750 R.P.M.?
    \item Un motor asincrónico trifásico de 1 par de polos y funciona con 60 Hz, si su velocidad nominal es de 3500 R.P.M ¿Cuál es su deslizamiento?
    \item El deslizamiento de un motor monofásico de 2 pares de polos es de 8\%. Calcular la velocidad de sincronismo y la velocidad del motor.
    \item El deslizamiento de un motor es de 10\% y su velocidad 1600 R.P.M. si el motor tiene 2 pares de polos ¿Qué frecuencia lo está alimentando?
    \item Calcular la velocidad de sincronismo de un motor que tiene una velocidad de 2801 R.P.M. y un deslizamiento del 6\%.
    \item Calcular el deslizamiento de un motor de 8 polos que gira a 800 R.P.M. y funciona a 60 Hz.
    \item Con que frecuencia será necesario alimentar un motor asincrónico trifásico de 1 HP para que tenga una velocidad de 1000 R.P.M si posee 2 pares de polos.
    \item El motor asincrono de 1.5HP anda a 1400 R.P.M ¿Cuál es el deslizamiento si es de 50 Hz y cuantos pares de polos tiene?
    \item El motor trifasico \emph{Marelli} de 4HP gira a 2890 R.P.M ¿Cuál es el deslizamiento si es de 50 Hz y cuantos pares de polos tiene?
    \item El motor monofásico asincrónico \emph{Industria Argentina} de 1450 R.P.M de 1HP ¿Qué deslizamiento posee y cuántos pares de polos tiene?
    \item Si los motores asincrónicos monofásicos normalizados tienen 1400 R.P.M ¿Cuál es el deslizamiento para un motor normalizado?
    \item En la plaza de un motor asincrónico monofásico de \emph{Industria Argentina} figuran los siguientes datos:
    \begin{itemize}
        \item 1HP
        \item 1425 R.P.M
        \item 50Hz
        \item 5.7A
        \item 220V
        \item Cap 430/514$\mu$
        \item $cos(\varphi)=0,9$
        \item IP21 uso general.
    \end{itemize}
         ¿Cuántos pares de polos tiene y cual es su rendimiento?
        \item En una publicación se vende un motor trifásico de velocidad 3000 R.P.M; marca \emph{WEG} si posee un 5\% de deslizamiento ¿Cuál es la velocidad del motor a carga nominal?
        \item Se determinó que un motor da 24 giros en 1 segundo, además el motor se alimenta con 220V 50Hz ¿Cuántos polos tiene el motor y cual es su deslizamiento?
        \item El eje de un motor tarda aproximadamente 20,8 milisegundos en dar una vuelta. Determinar si se alimenta con 50 Hz, el número de pares de polos, la velocidad de sincronismo y la velocidad del motor en R.P.M.
        \item Si un motor tiene un deslizamiento del 8\% y 2 pares de polos alimentando por 220V , 50Hz. Determinar en cuántos milisegundos da la vuelta el eje.
\end{enumerate}
\end{document}